\section{The trace of the complex Hessian}\label{sec:hessian}

The variations constructed in Section~\ref{var:section} let us apply
the second-derivative test at a maximizing disc. We compute the
central, logarithmically weighted, and unweighted terms separately,
then combine them in one inequality used by both positivity arguments.

Let $f$ be a maximizing disc furnished by Proposition~\ref{ex:maximizer}.
Write $x=f(0)$ and retain the K\"ahler metric $h$ used in its functional.
Use Proposition~\ref{var:frame} to choose a boundary-unitary holomorphic
frame $B=(b_1,\ldots,b_n)$ for $f^*T^{1,0}X$, and
Proposition~\ref{var:spray} to realize the sections $zb_j$ by a family
$F(t,z)$ with fixed center. The arguments in this section apply more
generally to any disc and family with those properties. After decreasing
the Sobolev exponent if necessary, all the maps and frames have
$W^{1,p_0}$ regularity for one $p_0>2$.

Set
\[
 N(z)=B(z)^*h(f(z))B(z).
\]
This positive-definite matrix is smooth in the open disc, continuous
and uniformly positive definite on its closure, and belongs to
$W^{1,p_0}$. Its boundary value is $I$.

\begin{lemma}[Determinant identity]\label{hes:determinant}
With these conventions,
\begin{equation}\label{eq:determinant-identity}
 \log\det N(0)=P_{R_h}(f).
\end{equation}
\end{lemma}

\begin{proof}
In a holomorphic target chart, the matrix of $B$ is holomorphic and
nonsingular. Thus
\[
 \partial_z\partial_{\bar z}\log\det N
 =\partial_z\partial_{\bar z}\log\det h(f(z))
 =-R_h(f',f').
\]
The first equality follows because $\log|\det B|^2$ is harmonic; the
formula is independent of the target chart. The right-hand side lies
in $L^{p_0/2}$. Apply Lemma~\ref{pre:green} and use
$\log\det N=0$ on the boundary.
\end{proof}

For a functional of the family, define
\[
 \cL=\left.\sum_{j=1}^n
   \frac{\partial^2}{\partial t_j\partial\overline{t_j}}\right|_{t=0}.
\]

\begin{proposition}[Trace identities]\label{hes:trace}
For every smooth closed real $(1,1)$-form $q$,
\begin{align}
 \cL d(F_t)&=\tr(B(0)^*k(x)B(0)),\label{eq:hessian-central}\\
 \cL P_q(F_t)&=\frac1{2\pi}\int_0^{2\pi}
      (\tr_h q)(f(e^{i\theta}))\,\dd\theta,\label{eq:hessian-poisson}\\
 \cL A_h(F_t)&=n-A_{R_h}(f).\label{eq:hessian-area}
\end{align}
Here $F_t$ means the disc $z\mapsto F(t,z)$, not a parameter derivative.
\end{proposition}

\begin{proof}
\textbf{The central term and differentiation under the integrals.}
Since $F(t,0)=x$, the vector $\partial_z F(t,0)$ is holomorphic in $t$
and lies in the fixed vector space $T_x^{1,0}X$. Its $t_j$ derivative
at zero is $b_j(0)$. Differentiating its squared $k(x)$ norm gives
Equation~\eqref{eq:hessian-central}. The holomorphic dependence of this
central derivative also follows directly from Cauchy's formula on a
smaller source circle.

We next justify the variations under the area and Poisson integrals.
Choose finitely many fixed source patches and target charts containing
the image of $F$ for sufficiently small $t$. In each chart the map
$t\mapsto F(t,\cdot)$ is holomorphic into $W^{1,p_0}$. The smooth metric
coefficients and the product rule imply that
\[
 t\longmapsto q(F_z(t,\cdot),F_z(t,\cdot))
\]
is twice continuously differentiable into $L^{p_0/2}$; the same holds
with $h$. Indeed, each parameter derivative is a sum of products of two
first source derivatives in $L^{p_0}$, with uniformly bounded Sobolev
coefficient factors. A fixed partition of unity in the source gives the
global assertion. Since $p_0/2>1$, integration against $1$ or against
$\log(1/|z|)$ is a continuous functional on this space. Differentiation
under both integrals is therefore valid.

\paragraph{The local identity and its Poisson integral.}
Locally write $q=\ddc\phi$. Joint holomorphy of $F$ on the open product
of the parameter domain and the disc gives
\begin{align*}
 q(F_z,F_z)&=\partial_z\partial_{\bar z}(\phi\circ F),\\
 \cL(\phi\circ F)&=\sum_j q(F_{t_j},F_{t_j})\big|_{t=0}.
\end{align*}
Commuting these derivatives and using $F_{t_j}(0,z)=zb_j(z)$ yields the
intrinsic identity
\begin{equation}\label{eq:local-hessian}
 \cL q(F_z,F_z)
 =\partial_z\partial_{\bar z}
     \left(|z|^2\tr(B^*q(f)B)\right).
\end{equation}
In particular, the mixed source derivative on the right belongs to
$L^{p_0/2}$ by the preceding differentiability argument.

The scalar function in parentheses is continuous on $\overline\D$ and
vanishes at zero. Lemma~\ref{pre:green} therefore computes its Poisson
integral from its boundary values. On the boundary, the columns of $B$
are an $h$-orthonormal basis, so
$\tr(B^*qB)=\tr_h q$. This proves
Equation~\eqref{eq:hessian-poisson}.

\paragraph{The area integral and its weak boundary limit.}
At the identity matrix, trace and log determinant have the same
differential. The boundary normalization will therefore let us replace
one radial derivative by the other, with an error controlled by the
radial Sobolev estimate along suitable radii.
For the area identity put $T=\tr N$. Integrating
Equation~\eqref{eq:local-hessian} with $q=h$ on $|z|<r$ gives
\begin{equation}\label{eq:area-radius-hessian}
 \frac1\pi\int_{|z|<r}\cL h(F_z,F_z)\,\dd S
 =r^2\langle T\rangle_r+\frac{r^3}{2}\langle T_r\rangle_r.
\end{equation}
The left-hand side tends to $\cL A_h(F_t)$ by integrability. The first
term on the right tends to $n$. To handle the second term without
assuming boundary differentiability, observe that
\begin{equation}\label{eq:trace-logdet-error}
 T_r-(\log\det N)_r=\tr((I-N^{-1})N_r).
\end{equation}
All matrix norms here are equivalent, and uniform positive definiteness
gives
\[
 \big|\langle T_r-(\log\det N)_r\rangle_r\big|
 \le C\|N(r,\cdot)-I\|_{L^2(S^1)}
       \|N_r(r,\cdot)\|_{L^2(S^1)}.
\]
The radial $W^{1,2}$ estimate and the trace $N=I$ give
\begin{equation}\label{eq:radial-trace-estimate}
 \|N(r,\cdot)-I\|_2
 \le(1-r)^{1/2}
       \left(\int_r^1\|N_s(s,\cdot)\|_2^2\,\dd s\right)^{1/2}
 =o((1-r)^{1/2}).
\end{equation}
This follows first on almost every radial line by the fundamental
theorem for Sobolev functions, and then by Cauchy--Schwarz and
integration over the angle. Smoothness inside the disc and continuity
at the boundary give the same estimate for the interior radii used
below.

Since $\int_{1/2}^1\|N_r\|_2^2\,\dd r<\infty$, there are radii
$r_\nu\uparrow1$ for which
$(1-r_\nu)\|N_r(r_\nu,\cdot)\|_2^2$ is bounded. Otherwise its eventual
lower bound would contradict integrability. Along these radii,
Equations~\eqref{eq:trace-logdet-error} and
\eqref{eq:radial-trace-estimate} show that the two circular derivative
means differ by $o(1)$.

On the other hand, the curvature calculation in
Lemma~\ref{hes:determinant} and Equation~\eqref{eq:green-area} give, for
every interior radius,
\[
 \frac r2\langle(\log\det N)_r\rangle_r=-A_{R_h,r}(f).
\]
The right-hand side tends to $-A_{R_h}(f)$. Substitute into
Equation~\eqref{eq:area-radius-hessian} along the chosen radii. Its
left-hand side has a limit independent of the sequence, proving
Equation~\eqref{eq:hessian-area}.
\end{proof}

\begin{corollary}[Maximizing inequality]\label{hes:maximum}
For a maximizer of $d+P_q-\delta A_h$ centered at $x$, the
boundary-unitary frame satisfies
\begin{equation}\label{eq:maximum-inequality}
 \tr(B(0)^*k(x)B(0))
 +\frac1{2\pi}\int_0^{2\pi}(\tr_hq)(f(e^{i\theta}))\,\dd\theta
 \le\delta\bigl(n-A_{R_h}(f)\bigr).
\end{equation}
It also satisfies Equation~\eqref{eq:determinant-identity}, and
$A_q(f)\ge-d(f)$.
\end{corollary}

\begin{proof}
Every sufficiently small parameter value gives an admissible
finite-area disc with the same center. Hence the real-valued twice
differentiable function $J(F_t)$ has a local maximum at zero, and
$\cL J(F_t)\le0$. Proposition~\ref{hes:trace} gives the displayed
inequality. The remaining assertions are Lemma~\ref{hes:determinant}
and the dilation inequality of Lemma~\ref{ex:dilation}.
\end{proof}

The two applications use different parts of the maximizing property.
In Section~\ref{sec:nefness}, we set $h=k$ and $q=R_k$ and let
$\delta$ tend to zero after comparing the maximum with each fixed test
disc; this gives a bound on $P_{R_k}$ for all finite-area discs.
In Section~\ref{sec:volume}, we set $q=-h$ and $\delta=1$.
Comparison with the constant disc then bounds $P_h(f)+A_h(f)$ by
$D$, independently of $h$. Together with $R_h\ge-h$, this is the
additional information that makes the volume estimate uniform over
all such metrics.
